\documentclass[11pt]{article}
\usepackage[T1]{fontenc}
\usepackage[utf8]{inputenc}
\usepackage{lmodern}
\usepackage[margin=1in]{geometry}
\usepackage{amsmath,amssymb,amsthm}
\usepackage{booktabs}
\usepackage{microtype}
\usepackage[hidelinks]{hyperref}

\newtheorem{theorem}{Theorem}
\newtheorem{lemma}{Lemma}
\newcommand{\F}{\mathcal F}
\newcommand{\Lz}{\mathcal L_z}
\newcommand{\Lzc}{\mathcal L_{z_C}}
\newcommand{\E}{\mathbb E}
\newcommand{\ind}{\mathbf 1}
\newcommand{\orb}{\operatorname{orb}}
\newcommand{\comp}{\text{compatible}}
\newcommand{\ip}[2]{\langle #1,#2\rangle}
\DeclareMathOperator{\tr}{tr}
\allowdisplaybreaks[1]
\setlength{\emergencystretch}{2em}

\title{New upper bounds for \texorpdfstring{$A(18,6,7)$, $A(19,6,7)$,\protect\\ and $A(19,6,8)$}{A(18,6,7), A(19,6,7), and A(19,6,8)}}
\author{Jiang Chang\thanks{ORCID: \href{https://orcid.org/0009-0006-8310-632X}{0009-0006-8310-632X}.}\\[3pt]
\small College of Computer Science and Software Engineering,\\
\small Hohai University, Nanjing, China\\
\small\href{mailto:science614@163.com}{science614@163.com}}
\date{}

\begin{document}
\maketitle

\begin{abstract}
We prove $A(18,6,7)\le309$, $A(19,6,7)\le460$, and
$A(19,6,8)\le667$, improving the previous upper bounds
$312$, $463$, and $673$, respectively. We augment Polak's four-point
semidefinite bound with anticode vacancy, integer occupancy, and cardinality
constraints obtained by established moment methods, and certify the resulting
bounds by exact rational dual certificates.
\end{abstract}

\section{Introduction}

Let $A(n,d,w)$ denote the largest cardinality of a binary code of length $n$
whose words have weight $w$ and whose distinct words have Hamming distance at
least $d$. Brouwer's tables \cite{Brouwer} collect bounds for this function.
For $(n,d,w)=(18,6,7),(19,6,7),(19,6,8)$, the previous integer upper bounds
are $312$, $463$, and $673$, respectively. The first two are due to
Polak \cite{Polak}; the third was established by Chang \cite{Chang} using an
exact certificate for Polak's four-point bound.

Semidefinite programming bounds for codes include Schrijver's three-point
bound based on the Terwilliger algebra \cite{Schrijver}, the quadruple-distance
bound of Gijswijt, Mittelmann, and Schrijver \cite{GMS}, and Polak's four-point
bounds for constant-weight codes \cite{Polak}. Strengthening an established
bound by additional valid inequalities also has precedents: Agrell, Vardy,
and Zeger \cite{AVZ} augmented linear programming inequalities for
constant-weight codes, and Kim and Toan \cite{KimToan} added linear constraints
to Schrijver's semidefinite program.

Here eight families of valid constraints strengthen Polak's bound at the
three parameters above. The constructions use Boolean multipliers and
linearization in the Sherali--Adams approach \cite{SheraliAdams}, localizing
matrices \cite{Lasserre}, finite-support and binomial moment inequalities
\cite{Prekopa}, symmetry reduction \cite{GatermannParrilo,Polak}, and classical
intersecting families as anticodes \cite{AhlswedeKhachatrian}. These are
established techniques; their application here gives the three improved
bounds and the accompanying exact certificates.

\section{The four-point bound and normalization}\label{sec:framework}

Write $[n]=\{1,\ldots,n\}$ and identify a word with its support in
\[
 \F=\binom{[n]}w,\qquad V=|\F|=\binom nw,\qquad
 d_H(u,v)=2\bigl(w-|u\cap v|\bigr).
\]
The relations indexed by $w-|u\cap v|$ form the Johnson scheme. Put
$\delta=\lceil d/2\rceil$. A set of words is \emph{compatible} if its distinct
members have mutual distance at least $d$.

For a code $C$, put $X_u(C)=\ind_{\{u\in C\}}$ and
$X_S=\prod_{u\in S}X_u$, with $X_\varnothing=1$. Boolean reduction uses
$X_u^2=X_u$ and $X_uX_v=0$ for incompatible distinct words. Thus
$X_SX_T=X_{S\cup T}$, with incompatible unions interpreted as zero.
For a uniformly distributed coordinate permutation $g\in S_n$, define
\begin{equation}\label{eq:normalization}
 z_C(S)=V\E_g[X_S(gC)],\qquad |S|\le4.
\end{equation}
These moments depend only on the orbit of the unordered set $S$. They satisfy
\[
 z_C(\varnothing)=V,\qquad z_C(\{u\})=|C|,\qquad
 0\le z_C(S)\le V.
\]
Write $z_1$ for the common singleton moment. For a Boolean-reduced polynomial
of degree at most four, set
\[
 \Lz\Bigl(\sum_S a_SX_S\Bigr)=\sum_S a_Sz(S).
\]
In particular, $\Lz(1)=V$ and $\Lzc(p)=V\E_g[p(gC)]$.
The affine constant of the semidefinite program is $1$; throughout the paper
$z(\varnothing)$ denotes the fixed number $V$. A term $az(\varnothing)$
therefore contributes the affine constant $aV$.

For an anchor $D\subseteq\F$, define
\[
 \mathcal I_D=\{S\subseteq\F:D\subseteq S,\ |D|+2|S\setminus D|\le4\},
 \qquad M_D(z)_{S,T}=z(S\cup T).
\]
The underlying four-point program has the constraints
\begin{equation}\label{eq:polak}
 \begin{split}
 M_\varnothing(z)&\succeq0,\qquad
 M_D(z)\succeq0\quad (D\text{ compatible},\ |D|=2),\\
 z(S)&\ge0\quad (S\text{ compatible},\ 1\le|S|\le4),
 \end{split}
\end{equation}
together with $z(\varnothing)=V$, orbit invariance, and zero moments for
incompatible sets. Maximizing $z_1$ gives Polak's bound $A_4(n,d,w)$
\cite{Polak}. In his normalization, $x(S)=z(S)/V$, so
$x(\varnothing)=1$ and $\sum_{u\in\F}x(\{u\})=z_1$.
The original constraints and their upper-bound property are those of
\cite[Section~1.1]{Polak}; its Sections~1.2 and~3 give the coordinate averaging
and representation block reductions. Section references to \cite{Polak}
refer to the linked arXiv version.

All additional constraints use the following positivity principle. If
$p(C)\ge0$ for every code, then $\Lzc(p)\ge0$. If $h(C)\ge0$ and $f$ is a real
vector of polynomials, then
\begin{equation}\label{eq:localization}
 \Lzc(hff^{\mathsf T})\succeq0,
 \qquad
 a^{\mathsf T}\Lzc(hff^{\mathsf T})a
 =V\E_g\bigl[h(gC)(a^{\mathsf T}f(gC))^2\bigr]\ge0,
\end{equation}
provided every reduced entry has degree at most four. Fixed congruences
$K\mapsto P^{\mathsf T}KP$, coordinate averaging, and division by positive
common factors preserve positivity. The added matrix families use retained
projections, as specified below; their exact rows and blocks are recorded
in the released models.

\section{Eight additional constraint families}\label{sec:constraints}

In the last three families, $B$ denotes an integer upper bound on the size of
every code with the specified parameters. Section~\ref{sec:baseline}
establishes this premise independently from the unaugmented four-point
program.

\subsection{Anticode vacancy scalars}\label{subsec:vacancy-scalar}

Put $t=w-\delta+1$. For $t\ge0$, $0\le r\le\delta-1$, and a coordinate set
$R$ of feasible size $t+2r$, let
\[
 Q(R,r)=\{u\in\F:|u\cap R|\ge t+r\},\qquad
 H_Q=1-\sum_{q\in Q}X_q.
\]
These are the classical intersecting families in
\cite[eq.~(1.9)]{AhlswedeKhachatrian}. For the present parameters,
$\delta=3$ and $r\in\{0,1,2\}$. For each retained compatible three-word
anchor $D$, the added scalar inequality is
\begin{equation}\label{eq:vacancy-scalar}
 z(D)-\sum_{q\in Q}z(D\cup\{q\})\ge0.
\end{equation}

\begin{proof}
For $u,v\in Q$,
$|u\cap v|\ge|u\cap R|+|v\cap R|-|R|\ge t$, so
$d_H(u,v)\le2(\delta-1)<d$. Hence $|C\cap Q|\le1$ and $H_Q(C)\in\{0,1\}$
for every code. The left side of \eqref{eq:vacancy-scalar} at $z_C$ is
$\Lzc(X_DH_Q)\ge0$ and uses at most four distinct words.
\end{proof}

The retained scalar rows have $D\cap Q=\varnothing$ and at least two
compatible extension words in $Q$. Duplicate rows are removed, and positions
of $R$ relative to $D$ are taken up to coordinate permutations.

\subsection{Anticode vacancy matrices}\label{subsec:vacancy-matrix}

Fix $Q=Q(R,r)$ and a word $a$. Let
\[
 \mathcal B_a=\{b\in\F:b=a\text{ or }d_H(a,b)\ge d\},\qquad
 K^{a,Q}(z)_{b,c}=z(\{a,b,c\})-\sum_{q\in Q}z(\{a,b,c,q\}).
\]
The constraints are the retained projections
\begin{equation}\label{eq:vacancy-matrix}
 P^{\mathsf T}K^{a,Q}(z)P\succeq0.
\end{equation}
Only harmonic degree zero is used here. More explicitly, the stabilizer of
$(a,R)$ permutes the four coordinate cells determined by membership in $a$
and $R$. If $\mathcal O_\beta$ is a word orbit, specified by its four cell
weights, the projection has columns
$P_{b,\beta}=\ind_{\{b\in\mathcal O_\beta\}}$. The projected entries are
$\sum_{b\in\mathcal O_\beta,c\in\mathcal O_\gamma}K^{a,Q}(z)_{b,c}$.
Degree-zero Young representatives give the same projection up to fixed
nonzero column normalizations.

\begin{proof}
For $v=(X_b)_{b\in\mathcal B_a}$,
$K^{a,Q}(z_C)=V\E_g[X_aH_Qvv^{\mathsf T}]\succeq0$ by
\eqref{eq:localization}. Its entries have degree at most four, and fixed
congruence gives \eqref{eq:vacancy-matrix}. The words $b,c$ are each compatible
with $a$; their entry is zero when they are incompatible with one another.
The index $b=a$ supplies the constant coordinate after multiplication by
$X_a$. Rows in $Q$ vanish under $H_Q$, and the whole matrix vanishes if
$a\in Q$; these zero rows and matrices are removed.
\end{proof}

\subsection{Matching-region occupancy projectors}\label{subsec:matching}

Choose $L\subseteq U\subseteq[n]$ and put
\[
 Q=\{u\in\F:L\subseteq u\subseteq U\},\qquad
 m=|U\setminus L|,\qquad k=w-|L|,
\]
with $\min(k,m-k)=\delta$ and $3\delta\le m<4\delta$.
Write $e_0=1$ and
\[
 e_\ell(X_Q)=\sum_{\substack{E\subseteq Q,\ |E|=\ell\\E\ \comp}}X_E,
 \qquad
 P_j=\sum_{\ell=j}^3(-1)^{\ell-j}\binom\ell j e_\ell(X_Q).
\]
Thus $P_0=1-e_1+e_2-e_3$, $P_1=e_1-2e_2+3e_3$,
$P_2=e_2-3e_3$, and $P_3=e_3$. For a word anchor $a$, the retained
constraints are
\begin{equation}\label{eq:matching}
 \sum_{\ell=j}^3(-1)^{\ell-j}\binom\ell j
 \sum_{\substack{E\subseteq Q,\ |E|=\ell\\E\ \comp}}
 z(\{a\}\cup E)\ge0,\qquad j=0,1,2.
\end{equation}

\begin{proof}
Deleting $L$ leaves $k$-subsets of an $m$-set. If $k=\delta$, the distance
condition makes them disjoint; if $m-k=\delta$, their complements are
disjoint. Thus $N=|C\cap Q|\le\lfloor m/\delta\rfloor=3$, and
$e_\ell(X_Q)(C)=\binom N\ell$. For $N\ge j$, binomial inversion gives
\[
 P_j(C)=\binom Nj\sum_{r=0}^{N-j}(-1)^r\binom{N-j}r
       =\ind_{\{N=j\}},
\]
and the same identity holds when $N<j$, with both sides zero.
Consequently \eqref{eq:matching} is $\Lzc(X_aP_j)\ge0$.
If $a\in Q$, the summation still allows $a\in E$, and set union merges its
repeated occurrence. Every term uses at most four distinct words.
\end{proof}

The $j=3$ inequality is a sum of nonnegative moments. The retained family
omits it, zero rows, duplicates, and other rows having only nonnegative
coefficients, and takes region and anchor representatives under coordinate
permutations.

\subsection{Coordinate interval spectra}\label{subsec:spectra}

First define an integer shortening bound, for $0\le k\le m$, by
\begin{equation}\label{eq:J}
 J(m,d,k)=
 \begin{cases}
 1,&\min(k,m-k)<\delta,\\[2pt]
 \displaystyle\min\left\{
 \left\lfloor\frac{mJ(m-1,d,k-1)}k\right\rfloor,
 \left\lfloor\frac{mJ(m-1,d,k)}{m-k}\right\rfloor\right\},&\text{otherwise}.
 \end{cases}
\end{equation}

\begin{lemma}\label{lem:capacity}
Every length-$m$, weight-$k$ code of minimum distance at least $d$ has
cardinality at most $J(m,d,k)$.
\end{lemma}
\begin{proof}
In the base case, two words have distance at most
$2\min(k,m-k)<d$. Otherwise, for a code $E$, let $E_i^1$ and $E_i^0$ denote
the words with a one or zero at coordinate $i$. Deleting that fixed
coordinate preserves distances. Induction and incidence counting give
\[
 k|E|=\sum_i|E_i^1|\le mJ(m-1,d,k-1),\qquad
 (m-k)|E|=\sum_i|E_i^0|\le mJ(m-1,d,k).
\]
Taking integer parts proves \eqref{eq:J}. The cases $k=0,m$ are base cases.
\end{proof}

For disjoint coordinate sets $L,Z$ with $|L|=\ell$ and $|Z|=b$, define
\[
 Q(L,Z)=\{u\in\F:L\subseteq u,\ u\cap Z=\varnothing\},\qquad
 N_{L,Z}=\sum_{u\in Q(L,Z)}X_u.
\]
The occupancy on a code is an integer in $\{0,\ldots,c\}$, where
$c=J(n-\ell-b,d,w-\ell)$. Only feasible intervals with $2\le c\le20$ are
selected. For a selected degree $1\le q\le\min(4,c)$, choose a $q$-element
set $\mathcal R\subseteq\{0,\ldots,c\}$ and $s\in\{1,-1\}$ such that
\begin{equation}\label{eq:spectral-polynomial}
 p(x)=s\prod_{r\in\mathcal R}(x-r)\ge0\quad(x=0,\ldots,c),\qquad
 p(x)=\sum_{r=0}^q a_r\binom xr,\quad a_r=\Delta^rp(0),
\end{equation}
where $\Delta p(x)=p(x+1)-p(x)$. The coefficients $a_r$ are integers.
The finite-support polynomial description is justified in
Appendix~\ref{app:spectra}.

Let $\Omega_r$ be a set of representatives of the orbits of compatible
unordered $r$-sets of words, and write
\[
 o(T)=|\orb(T)|,\qquad I(T)=\left|\bigcap_{u\in T}u\right|,\qquad
 U(T)=\left|\bigcup_{u\in T}u\right|,\qquad
 H_{\ell,b}=\binom n\ell\binom{n-\ell}b.
\]
The retained constraints are
\begin{equation}\label{eq:spectra}
 a_0H_{\ell,b}V+
 \sum_{r=1}^q a_r\sum_{T\in\Omega_r}
 o(T)\binom{I(T)}\ell\binom{n-U(T)}b\,z(T)\ge0.
\end{equation}
An impossible binomial coefficient is zero.

\begin{proof}
For each $(L,Z)$, Lemma~\ref{lem:capacity} gives $p(N_{L,Z}(C))\ge0$.
Boolean reduction represents $\binom{N_{L,Z}}r$ as the sum of $X_T$ over
compatible $r$-subsets of $Q(L,Z)$. Apply $\Lzc$ and sum over all
$H_{\ell,b}$ pairs. A nonempty $T$ lies in an interval exactly when $L$ lies
in its common intersection and $Z$ lies outside its union. This gives
$\binom{I(T)}\ell\binom{n-U(T)}b$ choices. Grouping by orbit proves
\eqref{eq:spectra}, including the constant $a_0H_{\ell,b}V$.
Its moments have order at most $q\le4$.
\end{proof}

\subsection{Coordinate occupancy matrices}\label{subsec:occupancy}

Use either orientation $\varepsilon\in\{0,1\}$, with
\[
 u^{(0)}=u,\qquad u^{(1)}=[n]\setminus u,\qquad w_0=w,\quad w_1=n-w.
\]
Complementation preserves distance, and the indicator of $u^{(1)}$ in the
complemented code is $X_u$. For a nonempty core $R$, $|R|=\ell\le w_\varepsilon$,
put
\[
 N_R=\sum_{u:R\subseteq u^{(\varepsilon)}}X_u,\qquad
 F_{R,i}=\sum_{u:|u^{(\varepsilon)}\cap R|=i}X_u,\qquad
 v_R=(1,(F_{R,i})_i)^{\mathsf T},
\]
where $\max(0,\ell+w_\varepsilon-n)\le i\le\min(\ell,w_\varepsilon)$.
Lemma~\ref{lem:capacity} gives $0\le N_R(C)\le c$, with
$c=J(n-\ell,d,w_\varepsilon-\ell)$. The selected pairs $(\ell,c)$ are listed
in Table~\ref{tab:cores}. For each pair, use
\[
 p_j(x)=(x-j)(x-j-1)\quad(0\le j<c),\qquad
 p_{\mathrm{end}}(x)=x(c-x),
\]
and impose
\begin{equation}\label{eq:occupancy}
 \mathcal O_{\varepsilon,\ell,p}(z)
 =\sum_{\substack{R\subseteq[n]\\|R|=\ell}}
 \Lz\bigl(p(N_R)v_Rv_R^{\mathsf T}\bigr)\succeq0.
\end{equation}

\begin{table}[htbp]
\centering
\caption{Selected core sizes and capacities for coordinate occupancy matrices.}
\label{tab:cores}
\small
\begin{tabular}{@{}llll@{}}
\toprule
$(n,d,w)$ & Original orientation & Complement orientation & Matrices\\
\midrule
$(18,6,7)$ & $(3,15),(4,4)$ & $(6,19),(7,8),(8,3)$ & $54$\\
$(19,6,7)$ & $(3,20),(4,5)$ & $(7,19),(8,8),(9,3)$ & $60$\\
$(19,6,8)$ & $(4,15),(5,4)$ & $(7,9),(8,3)$ & $35$\\
\bottomrule
\end{tabular}
\end{table}

\begin{proof}
For an integer $N_R$, the consecutive factors $N_R-j$ and $N_R-j-1$
have nonnegative product. The capacity gives $N_R(c-N_R)\ge0$.
Every summand in \eqref{eq:occupancy} is therefore positive semidefinite at
$z_C$ by \eqref{eq:localization}, and its reduced degree is at most
$2+1+1=4$. Coordinate invariance gives the exact normalization
\begin{equation}\label{eq:occupancy-normalization}
 \begin{split}
 \mathcal O_{\varepsilon,\ell,p}(z_C)
 &=V\sum_{|R|=\ell}p(N_R(C))v_R(C)v_R(C)^{\mathsf T}\\
 &=V\binom n\ell\E_R[p(N_R(C))v_R(C)v_R(C)^{\mathsf T}].
 \end{split}
\end{equation}
Each selected pair contributes $c+1$ matrices. The vector $v_R$ is a fixed
linear image of the constant and the individual-word indicators, so this
construction is also a fixed congruence of a localizing matrix.
\end{proof}

Appendix~\ref{app:coefficients} gives the orbit coefficients, including all
repeated-word terms and both orientations. Zero-row deletion and removal of
positive common factors preserve the displayed constraints.

\subsection{Global cardinality cap}\label{subsec:cap}

Let $M(X)=\sum_{u\in\F}X_u$. The independently certified premise
$|C|\le B$ gives $B-M(X)\ge0$ on every code. The global constraint is
\begin{equation}\label{eq:cap}
 B-z_1\ge0.
\end{equation}
\begin{proof}
At $z_C$, the left side is $B-|C|\ge0$. Its normalization is also given by
$V^{-1}\Lz(B-M)=B-V^{-1}\sum_{u\in\F}z(\{u\})=B-z_1$.
Thus the affine constant in this row is $B$.
\end{proof}

\subsection{Anchored cardinality scalars}\label{subsec:cardinality-scalar}

For each retained nonempty compatible anchor $D$, $k=|D|\le3$, impose
\begin{equation}\label{eq:cardinality-scalar}
 Bz(D)-\sum_{u\in\F}z(D\cup\{u\})
 =(B-k)z(D)-\sum_{u\in\F\setminus D}z(D\cup\{u\})\ge0.
\end{equation}
\begin{proof}
The first expression at $z_C$ is $\Lzc(X_D(B-M))\ge0$.
The integrand is nonnegative for every code of size at most $B$, and its
reduced degree is at most four. Repeated anchor words give the $kz(D)$ term.
\end{proof}

For compatible orbit representatives $T\in\Omega_k$ and
$W\in\Omega_{k+1}$, let
\begin{equation}\label{eq:extensions}
 m(T,W)=\#\{T'\subseteq W:|T'|=k,\ T'\in\orb(T)\},\qquad
 e(T,W)=\frac{o(W)m(T,W)}{o(T)}.
\end{equation}
The exact orbit form of \eqref{eq:cardinality-scalar} is
\begin{equation}\label{eq:cardinality-orbit}
 (B-k)z(T)-\sum_{W\in\Omega_{k+1}}e(T,W)z(W)\ge0.
\end{equation}
The integer incidence count $e(T,W)$ and the treatment of repeated words
are proved in Appendix~\ref{app:extensions}.

\subsection{Rooted cardinality matrices}\label{subsec:cardinality-matrix}

For a word $a$ and $\mathcal B_a$ as above, define
\[
 K^{a,B}(z)_{b,c}
 =Bz(\{a,b,c\})-\sum_{u\in\F}z(\{a,b,c,u\}),
 \qquad b,c\in\mathcal B_a.
\]
The added constraints are the retained Young projections
\begin{equation}\label{eq:cardinality-matrix}
 P^{\mathsf T}K^{a,B}(z)P\succeq0.
\end{equation}

\begin{proof}
For $v=(X_b)_{b\in\mathcal B_a}$,
\[
 K^{a,B}(z_C)
 =V\E_g[X_a(B-M)vv^{\mathsf T}]
 =(B-|C|)\Lzc(X_avv^{\mathsf T})\succeq0.
\]
All entries have at most four distinct words; repetitions, including $b=a$
or $c=a$, are merged by union. Fixed congruence proves
\eqref{eq:cardinality-matrix} for every code.
\end{proof}

The stabilizer of $a$ is $S_w\times S_{n-w}$. Its projections use tensor
products of binary Young representatives of shapes $(w-h,h)$ and
$(n-w-j,j)$, where $h,j\ge0$, $2h\le w$, and $2j\le n-w$. The weight index
$s=|a\cap b|$ satisfies
\[
 h\le s\le w-h,\qquad j\le w-s\le n-w-j,\qquad
 s=w\ \text{or}\ w-s\ge\delta.
\]
This is the singleton-anchor Young construction of
\cite[Section~3.1]{Polak}, applied to the cardinality localizer.
For compatible $T=\{a,b,c\}$ of size $k$, the corresponding entry is exactly
the left side of \eqref{eq:cardinality-orbit}; incompatible entries are zero.
This linear substitution commutes with the fixed projections.

\section{The independently certified cardinality premise}\label{sec:baseline}

The three baseline certificates concern only the original Polak program
\eqref{eq:polak}. They give
\begin{align*}
 U_{\mathrm{base}}(18,6,7)
 &=\frac{13433029372173119186202113540456444959822875841641987}
 {42949672960000000000000000000000000000000000000000},\\[4pt]
 U_{\mathrm{base}}(19,6,7)
 &=\frac{1274550274376814049016394972588387247712471375324415947}
 {2748779069440000000000000000000000000000000000000000},\\[4pt]
 U_{\mathrm{base}}(19,6,8)
 &=\frac{841300017883100309254921891583162664400313}
 {1250000000000000000000000000000000000000}.
\end{align*}
Every code gives a feasible point of that program, so exact weak duality
gives $|C|\le U_{\mathrm{base}}$. The integrality of $|C|$ then gives
\[
 B=\lfloor U_{\mathrm{base}}\rfloor=
 \begin{cases}
 312,&(n,d,w)=(18,6,7),\\
 463,&(n,d,w)=(19,6,7),\\
 673,&(n,d,w)=(19,6,8).
 \end{cases}
\]
This establishes $B-M(X)\ge0$ on all codes, including the empty code and
codes of size smaller than $B$, and justifies the last three families in
Section~\ref{sec:constraints}. The third baseline bound is the certificate
bound of \cite{Chang}.

\section{Exact dual certificates and verification}\label{sec:certificates}

The proofs above show that every code induces a feasible point of the
augmented program, with objective $z_1=|C|$. Only moments of at most four
distinct words occur. We give the moment bounds needed to control
coefficient residuals on the entire semidefinite feasible region.

\begin{lemma}\label{lem:box}
Every feasible point of the invariant four-point program satisfies
$0\le z(S)\le V$ for $|S|\le4$.
\end{lemma}
\begin{proof}
Nonnegativity is imposed in \eqref{eq:polak}, and incompatible moments are
zero. The principal submatrix of $M_\varnothing$ on
$\varnothing,\{u\}$ gives
\[
 \begin{pmatrix}V&z_1\\z_1&z_1\end{pmatrix}\succeq0,
 \qquad z_1^2\le Vz_1,
\]
so $z_1\le V$. For a compatible pair, the submatrix on
$\{u\},\{u,v\}$ gives $z(\{u,v\})^2\le z_1z(\{u,v\})$,
hence $z(\{u,v\})\le z_1$. For a compatible triple, use the indices
$D,D\cup\{w\}$ in $M_D$, where $D=\{u,v\}$, to obtain
$z(D\cup\{w\})^2\le z(D)z(D\cup\{w\})$.
Finally, the entry between $D\cup\{w\}$ and $D\cup\{x\}$ gives
\[
 z(D\cup\{w,x\})^2
 \le z(D\cup\{w\})z(D\cup\{x\})\le V^2.
\]
These inequalities also cover zero diagonal entries and establish the claim.
\end{proof}

To state exact weak duality, index the nonempty orbit variables by $i$,
write the objective as $c^{\mathsf T}z=z_1$, and write every primal block as
$F_j(z)=F_{j,0}+\sum_iF_{j,i}z_i\succeq0$; a scalar inequality is a
$1\times1$ block. For rational positive semidefinite dual matrices $Y_j$, put
\[
 \beta=\sum_j\ip{Y_j}{F_{j,0}},\qquad
 r_i=c_i+\sum_j\ip{Y_j}{F_{j,i}},\qquad
 \ip{Y}{F}=\tr(YF).
\]
Exact arithmetic gives the identity
\begin{equation}\label{eq:weak-dual}
 z_1=\beta+\sum_i r_i z_i-\sum_j\ip{Y_j}{F_j(z)}
 \le\beta+\sum_i r_i z_i.
\end{equation}
Lemma~\ref{lem:box} bounds the residual term by
$V\sum_i\max(r_i,0)$. Thus a recorded rational residual allowance $\eta$
that bounds this term yields $z_1\le\beta+\eta$ for every feasible point.
The certified upper bounds include this allowance.

Each certificate records the rational dual matrices, exact
$LDL^{\mathsf T}$ factors, residual compensation, and hashes of the
corresponding JSON and SDPA models. Exact verification establishes positive
semidefiniteness and the coefficient identities and bounds in
\eqref{eq:weak-dual}. The same argument applies to the independent baseline
certificates and to the augmented certificates. The dual matrices were obtained by rounding numerical solutions of the semidefinite programs to rational matrices; the bounds rely only on the exact verification described above.

\section{Results}\label{sec:results}

\begin{theorem}
The exact rational upper bounds in Table~\ref{tab:results} hold. In particular,
\[
 \boxed{A(18,6,7)\le309,\qquad A(19,6,7)\le460,\qquad A(19,6,8)\le667.}
\]
\end{theorem}
\begin{proof}
For each parameter triple, the baseline certificate establishes the
cardinality premise of Section~\ref{sec:baseline}. Every code moment vector
then satisfies all eight added constraint families. Applying the augmented
certificate and \eqref{eq:weak-dual} bounds its objective $|C|$ by the stated
rational number. Taking the integer part gives the three assertions.
\end{proof}

\begin{table}[htbp]
\centering
\caption{Exact certified bounds. The decimal expansions are complete.}
\label{tab:results}
\footnotesize
\setlength{\tabcolsep}{4pt}
\renewcommand{\arraystretch}{1.75}
\begin{tabular}{@{}cccc@{}}
\toprule
$(n,d,w)$ & Rational upper bound & Exact decimal & Integer bound\\
\midrule
$(18,6,7)$ &
$\dfrac{3868791767977833258518625959723}{12500000000000000000000000000}$ &
$309.50334143822666068149007677784$ & $309$\\[4pt]
$(19,6,7)$ &
$\dfrac{46047385392564732804912812389921}{100000000000000000000000000000}$ &
$460.47385392564732804912812389921$ & $460$\\[4pt]
$(19,6,8)$ &
$\dfrac{16683417296961520597884426836631}{25000000000000000000000000000}$ &
$667.33669187846082391537707346524$ & $667$\\
\bottomrule
\end{tabular}
\end{table}

\section{Data and code availability}\label{sec:data}

The certificate package is available at
\begin{center}
\url{https://github.com/wille614/cw-code-upper-bounds}
\end{center}
as release \href{https://github.com/wille614/cw-code-upper-bounds/releases/tag/v1.0.0}{v1.0.0}.
It contains the three baseline and three augmented models, each in JSON
and SDPA form, their six rational dual certificates, an input preparation
program, and the exact verifier. The 18 compressed assets have published
SHA-256 checksums. The baseline certificate for $(19,6,8)$ is byte-identical
to the published certificate accompanying \cite{Chang}.

Offline verification uses CPython 3.11 or later, its standard library, and
enabled assertions. For each parameter triple, the preparation program
checks the six asset hashes, decompresses the models, and checks the
augmented model's reference to the baseline certificate, its rational bound,
and its integer part. Separate verifier invocations establish the baseline
premise and the augmented bound. The release documentation gives commands
for Windows, Linux, and macOS. The six verifications have integer-bound
pairs $(312,309)$, $(463,460)$, and $(673,667)$; successful runs report
\texttt{PASS} and exit code $0$.

For $(18,6,7)$, $(19,6,7)$, and $(19,6,8)$, respectively, the augmented
verifications check $46677$, $52716$, and $116007$ blocks, covering
$13834807$, $15981482$, and $66812277$ SDPA coefficients. The code is
licensed under Apache License 2.0, and the certificates and model data under
CC BY 4.0.

\section*{Acknowledgements}

The experiments and the writing were carried out with AI assistance.

\appendix

\section{Finite-support polynomials}\label{app:spectra}

For $q\le c$, let $\mathcal P_q$ be the cone of real polynomials of degree
at most $q$ that are nonnegative at $0,\ldots,c$. Evaluation at these points
is injective, so the cone is pointed; the strictly positive constant
polynomial shows that it has dimension $q+1$. An extreme ray must have at
least $q$ independent active evaluation constraints. A nonzero polynomial
of degree at most $q$ has at most $q$ distinct roots. Thus every extreme ray
is a positive multiple of a signed product in
\eqref{eq:spectral-polynomial}. Conversely, its $q$ distinct roots give
$q$ independent equations, so a valid signed product spans a one-dimensional
face. Lower-degree valid products are also allowed. The Newton expansion
$p(x)=\sum_r\Delta^rp(0)\binom xr$ gives the coefficients used in
\eqref{eq:spectra}.

\section{Orbit coefficients and repeated words}\label{app:coefficients}

\subsection{Coordinate occupancy matrices}

Expand an entry of $N_R^p v_Rv_R^{\mathsf T}$ for $p=0,1,2$. There are at
most four ordered word slots: $p$ region slots and zero, one, or two row
slots; the constant row contributes no slot. Repetitions are allowed.
Every tuple is specified uniquely by its set $T$ of distinct words and a
surjection from the slots onto $T$. Incompatible $T$ contribute zero.
This accounts for the diagonal terms in $N_R^2$ and for coincidences between
row and region words.

For a representative $T$ and a fixed surjection, the core $R$ must lie in
the intersection of all oriented region words. With no region slot this
eligible set is $[n]$. Partition the eligible set into four cells by the
two oriented row-word bits, using bit zero for a missing row. If the sizes
are $c_{00},c_{01},c_{10},c_{11}$, the number of possible cores is
\begin{equation}\label{eq:core-coefficient}
 \sum_{(x_{00},x_{01},x_{10},x_{11})}
 \prod_{a,b\in\{0,1\}}\binom{c_{ab}}{x_{ab}},
\end{equation}
where $\sum_{a,b}x_{ab}=\ell$ and, for present row slots, the intersection
requirements are $x_{10}+x_{11}=i$ and $x_{01}+x_{11}=j$.
Multiplication by $o(T)$, summation over the surjections, and replacement
of $X_T$ by $z(T)$ give the coefficient for that orbit. The empty tuple
contributes $\binom n\ell V$. Combining the powers $p=0,1,2$ with the
coefficients of the chosen quadratic gives exactly
\eqref{eq:occupancy}.

For the complement orientation, all coordinate bits are reversed before
these counts are formed. The moment is still that of the original $T$,
with the same orbit size. Thus both orientations use the same orbit
variables.

\subsection{Cardinality extension counts}\label{app:extensions}

To justify \eqref{eq:extensions}, count pairs $(T',W')$ of the prescribed
orbit types with $T'\subseteq W'$. There are $o(W)m(T,W)$ such incidences.
Transitivity makes the number above each $T'$ equal, so
$e(T,W)=o(W)m(T,W)/o(T)$ is an integer. For fixed $T$, each containing
$W'$ corresponds to exactly one compatible new word extending $T$.
Separating the $k=|T|$ repeated words in the sum over all $u\in\F$ gives
\eqref{eq:cardinality-orbit}.

For code-induced moments, invariance of $|gC|$ also gives
\[
 |C|z_C(T)=kz_C(T)+\sum_{W\in\Omega_{k+1}}e(T,W)z_C(W).
\]
For an arbitrary semidefinite vector, the expression on the right is the
linearized moment $\Lz(MX_T)$. Applying this expression entrywise to
$T=\{a,b,c\}$ gives the rooted cardinality matrix before its fixed Young
projection. Orbit averaging, deletion of zero or repeated rows, and
positive rescaling preserve the validity of all retained constraints.

\begin{thebibliography}{99}

\bibitem{Brouwer}
A.~E. Brouwer,
\emph{Bounds for binary constant weight codes}, online table,
\url{https://aeb.win.tue.nl/codes/Andw.html}, accessed 5 November 2026.

\bibitem{Schrijver}
A. Schrijver,
New code upper bounds from the Terwilliger algebra and semidefinite programming,
\emph{IEEE Trans. Inform. Theory} \textbf{51} (2005), no.~8, 2859--2866.
\href{https://doi.org/10.1109/TIT.2005.851748}{doi:10.1109/TIT.2005.851748}.

\bibitem{GMS}
D.~C. Gijswijt, H.~D. Mittelmann, and A. Schrijver,
Semidefinite code bounds based on quadruple distances,
\emph{IEEE Trans. Inform. Theory} \textbf{58} (2012), no.~5, 2697--2705.
\href{https://doi.org/10.1109/TIT.2012.2184845}{doi:10.1109/TIT.2012.2184845}.

\bibitem{Polak}
S.~C. Polak,
Semidefinite programming bounds for constant-weight codes,
\emph{IEEE Trans. Inform. Theory} \textbf{65} (2019), no.~1, 28--38.
\href{https://doi.org/10.1109/TIT.2018.2854800}{doi:10.1109/TIT.2018.2854800}.
\href{https://arxiv.org/abs/1703.05171}{arXiv:1703.05171}.

\bibitem{AVZ}
E. Agrell, A. Vardy, and K. Zeger,
Upper bounds for constant-weight codes,
\emph{IEEE Trans. Inform. Theory} \textbf{46} (2000), no.~7, 2373--2395.
\href{https://doi.org/10.1109/18.887851}{doi:10.1109/18.887851}.

\bibitem{KimToan}
H.~K. Kim and P.~T. Toan,
Improved semidefinite programming bound on sizes of codes,
\emph{IEEE Trans. Inform. Theory} \textbf{59} (2013), no.~11, 7337--7345.
\href{https://doi.org/10.1109/TIT.2013.2277714}{doi:10.1109/TIT.2013.2277714}.
\href{https://arxiv.org/abs/1212.3467}{arXiv:1212.3467}.

\bibitem{SheraliAdams}
H.~D. Sherali and W.~P. Adams,
A hierarchy of relaxations between the continuous and convex hull
representations for zero-one programming problems,
\emph{SIAM J. Discrete Math.} \textbf{3} (1990), no.~3, 411--430.
\href{https://doi.org/10.1137/0403036}{doi:10.1137/0403036}.

\bibitem{Lasserre}
J.~B. Lasserre,
Global optimization with polynomials and the problem of moments,
\emph{SIAM J. Optim.} \textbf{11} (2001), no.~3, 796--817.
\href{https://doi.org/10.1137/S1052623400366802}{doi:10.1137/S1052623400366802}.

\bibitem{Prekopa}
A. Pr\'ekopa,
Boole--Bonferroni inequalities and linear programming,
\emph{Oper. Res.} \textbf{36} (1988), no.~1, 145--162.
\href{https://doi.org/10.1287/opre.36.1.145}{doi:10.1287/opre.36.1.145}.

\bibitem{GatermannParrilo}
K. Gatermann and P.~A. Parrilo,
Symmetry groups, semidefinite programs, and sums of squares,
\emph{J. Pure Appl. Algebra} \textbf{192} (2004), nos.~1--3, 95--128.
\href{https://doi.org/10.1016/j.jpaa.2003.12.011}{doi:10.1016/j.jpaa.2003.12.011}.

\bibitem{AhlswedeKhachatrian}
R. Ahlswede and L.~H. Khachatrian,
The complete intersection theorem for systems of finite sets,
\emph{European J. Combin.} \textbf{18} (1997), no.~2, 125--136.
\href{https://doi.org/10.1006/eujc.1995.0092}{doi:10.1006/eujc.1995.0092}.

\bibitem{Chang}
J. Chang,
A new upper bound for $A(19,6,8)$,
arXiv:2611.00173 (2026).
\url{https://arxiv.org/abs/2611.00173}.

\end{thebibliography}
\end{document}
